\subsection{Baselines}
\label{subsection:baselines}

We compare our best-performing \textsf{\small SB-CLASSIFIER} with an adapted version in which the URL classifier is replaced by an (unrealistic) perfect oracle~(\textsf{\small SB-ORACLE}), as well as with several baseline methods. First, we consider four simple crawlers.
\begin{inparaenum}[(i)]
\item The \textsf{\small RANDOM} crawler selects the next hyperlink to visit
	uniformly at random from the crawl frontier.
	\item The
    \emph{Breadth-First Search} (\textsf{\small BFS}) exhaustive crawler maintains the
	frontier as a \emph{First-In-First-Out queue} data structure. It crawls
	all pages reachable by a path of length $\ell$ before any page reachable
	only by longer paths ($\ell' > \ell$).
	\item \emph{Depth-First
    Search} (\textsf{\small{DFS}}) maintains the frontier in a \emph{Last-In-First-Out stack}.
	It is rarely used for exhaustive crawling since it may fall
        into
	\href{https://en.wikipedia.org/wiki/Spider_trap}{robot traps}.
	\item The unrealistic \emph{Omniscient Crawler} (\textsf{\small OMNISCIENT}) knows all target URLs ($V^*$) from the start and crawls them sequentially. Since an optimal crawler is intractable (Prop.~\ref{prop:np-complete}), this omniscient crawler provides an (unreachable) upper bound on crawler efficiency.
\end{inparaenum}


The \emph{Offline-Trained, Tag Path-Based Crawler} (\textsf{\small TP-OFF})~\cite{faheem2015adaptive} is a closely related focused crawler originally designed to retrieve \emph{diverse} textual content. We apply it \emph{technically unchanged}, except for redefining its reward to target retrieval.
As described in~\cite{faheem2015adaptive}, it first crawls 3k pages with BFS and groups the tag paths leading to followed links as in Sec.~\ref{subsection:states_actions}. Page benefits are computed immediately, and tag path groups are stored in a priority queue ordered by average benefit. After these 3k pages, \cite{faheem2015adaptive} \emph{only considers links matching existing tag path groups} (ordered by priority), other are ignored. 
In our context, benefit computation requires knowing if links lead to targets. To enable this baseline, \emph{we provide it with the true benefit} (number of targets behind a page's links) as if given by an oracle, for the initial 3k pages. Beyond this point, newly formed tag path groups receive a fixed benefit of~0. This baseline, given an unfair advantage in its first stage, can be seen as an \emph{ablated} version of ours, learning \emph{offline} instead of \emph{online} with our RL method.

The \emph{Focused Crawler} (\textsf{\small FOCUSED}) represents early generic focused crawling approaches that rely on a classifier to estimate the likelihood that a newly discovered hyperlink leads to a target~\cite{chakrabarti1999focused,diligenti2000focused}. The frontier is maintained as a \emph{priority queue},  favoring links predicted to be relevant. It relies on a standard \emph{logistic regression}, periodically retrained on crawled pages at no extra HTTP cost.
Its features follow standard focused crawler practice: the (approximated) depth of the source page, a 2-gram BoW representation of the URL (similarly to Sec.~\ref{subsection:url_classifier}), and a 2-gram BoW representation of the link's \emph{anchor} text~\cite{chakrabarti1999focused,diligenti2000focused}. \emph{Topic-oriented} features are intentionally excluded (further discussed in Sec.~\ref{section:related_work}). Overall, \textsf{\small FOCUSED} can be viewed as an \emph{ablated} version of our crawler, as it neither exploits tag-path structure (Sec.~\ref{subsection:url_classifier}) nor uses RL (Sec.~\ref{subsection:crawling_algorithm}).


{\textsf{\small TRES}}~\cite{kontogiannis2022tree} is a topical, RL-based focused crawler using a Bi-LSTM classifier to assess HTML page relevance. Like our approach, it employs RL to learn where interesting content resides in a website. However, {\textsf{\small TRES}} defines interestingness through a user-specified topic, whereas we target SDs regardless of subject. As it targets topic-relevant HTML pages only, it requires input keywords, positive HTML examples, and is limited to one language.
To include {\textsf{\small TRES}} as a baseline for SD retrieval, we made changes related to its focus, \emph{without altering its core logic or technical components}. Specifically, we give it three \emph{unfair advantages}: \begin{inparaenum}[(i)]
\item \textsf{\small TRES} expects a list of topic-specific keywords to initialize its classifier's training: we supply 74 hand-crafted terms likely to appear in links' anchors to targets to initialize its classifier.  
\item We provide 1k HTML pages with links to targets from the evaluation websites as positive examples (based on prior crawls), as required to pre-train its classifier, giving it partial access to the solution.  
\item \textsf{\small TRES} must classify URLs as HTML or not to manage its frontier, as it only accepts HTML pages. Although costly in practice (see Sec.~\ref{subsection:url_classifier}), we gave it an oracle doing this classification at no cost, giving it its chance while not affecting its original behavior.
\end{inparaenum}
Since its RL agent and classifier rely only on textual HTML features, {\textsf{\small TRES}} cannot directly detect targets. We therefore added two adjustments: links are added to the frontier as in~\cite{kontogiannis2022tree}, while others (which TRES would normally ignore) are visited immediately and targets counted if found; and its language filter was disabled, preventing coarse URL-based false exclusions. The modified code is available at~\cite{tres_modified_repository}.

\begin{toappendix}
  \subsection{Baselines}
    We use the following list of keywords as initial input to
  \textsf{\small TRES} to train its classifier:
    \begin{multicols}{4}
      \noindent
      pdf\\
xls\\
csv\\
tar\\
zip\\
rar\\
rdf\\
json\\
doc\\
xml\\
yaml\\
txt\\
tsv\\
ppt\\
ods\\
dta\\
7z\\
ttl\\
file\\
document\\
report\\
publication\\
dataset\\
data\\
download\\
archive\\
spreadsheet\\
table\\
list\\
resource\\
annex\\
supplement\\
attachment\\
proceedings\\
survey\\
material\\
output\\
content\\
statistics\\
article\\
paper\\
metadata\\
fact\\
download file\\
download document\\
available for download\\
access data\\
view report\\
get dataset\\
data file\\
read more\\
resource list\\
get document\\
download pulication\\
document archive\\
supporting materials\\
export data\\
download CSV\\
download PDF\\
download XLS\\
dataset download\\
attached document\\
official documents\\
browse files\\
download statistics\\
download article\\
annual report\\
white paper\\
technical documentation\\
technical report\\
raw data\\
metadata file\\
open data\\
fact sheet\\

    \end{multicols}
  
\end{toappendix}
